\origpage{379}{N. Robertson, P.D. Seymour / Journal of Combinatorial Theory, Series B 90 (2004) 325--385}
\noindent \del{union of}
\begin{equation*}
\delm{(G' \cap \text{sk}(\Gamma_j)) \cup \bigcup(P(e) : e \in E')}
\end{equation*}
\del{with the graph with vertex set $V(M(\Gamma_j, \mathscr{L}_j))$ and no edges.}

{\sloppy \del{For each $v \in V(\Gamma'_i)$ there is a component $K$ of $G$ with $V(\sigma'(v)) \subseteq V(K)$; because for each $e \in E(\sigma(v)) \setminus E(\text{sk}(\Gamma_j)) \subseteq E'$, there is a path of $P(e)$ joining the ends of $e$. Let $\sigma''(v) = K$. For $v \in V(M(\Gamma_i, \mathscr{L}_i)) \setminus V(\Gamma'_i)$, if $v \in V(H(W)) \cap V(\Gamma_i)$ for some $W \in \mathscr{W}$, let $\sigma''(v)$ be the component of $G$ with $v \in V(\sigma''(v))$. If $v$ is the $k\text{th}$ term of $\mu_i(u)$ for some $u \in V(\Gamma_i) \cap \text{bd}(W)$, let $\sigma''(v)$ be the component of $G$ containing the $k\text{th}$ term of $\mu_j(u')$, where if $e$ is the border edge of $\Gamma'_i$ with tail $u$ then $u'$ is the tail of $\sigma'(e)$. For each $v$, let $\sigma(v)$ be the induced subgraph of $\text{sk}(\Gamma)$ with vertex set $V(\sigma''(v))$. For each edge $e$ of $M(\Gamma_i, \mathscr{L}_i)$, let $\sigma(e) = \sigma'(e)$. We claim that $\sigma$ is an inflation of $M(\Gamma_i, \mathscr{L}_i)$ in $\Gamma_j$, as in the definition of outline. We omit the verification, because it is lengthy and almost identical with (but easier than) the major part of the proof of Theorem 9.1 of [11].}\par}

\fix{We give a different construction of $\sigma$, with its verification. Let $W \in \mathscr{W}$, choose $r_0^* \in S(W)$, and fix an orientation of $bd(W)$. Let $u_0, \dots, u_{k-1}$ be the vertices of $\Gamma_i$ in $bd(W)$ in order, where $r_0^*$ lies between $u_{k-1}$ and $u_0$, and let $w_0, \dots, w_{l-1}$ be those of $\Gamma_j$, numbered in the same way. Since the neighbours of $r_0^*$ are in $N(W)$, we have $w_0 = u_0$ and $w_{l-1} = u_{k-1}$. For $t \ge 0$ let $u_p^t$ be the $(t+1)$th term of $\mu_i(u_p)$, the vertex at \textit{level} $t$ of $S_{u_p}$, and define $w_q^t$ in the same way. In the definition of $\phi'_i(e_{r^*})$ we take $v_1$ to be the tail of $e_{r^*}$. Two spokes with the same inner end have the same number of levels, since each spoke has $\rho(W) - 1$ or $\rho(W)$ edges, by the last condition in the definition of a wheel, and $\Gamma^*$ is bipartite. They also have the same hub index. For if they have an even number of edges, then the last level of each is their common inner end; and otherwise the inner end is not in $V(\Gamma)$, and the last level of each is not in $V(H(W))$, since a spoke meets $V(H(W))$ only in its inner end, and so both hub indices are $0$.

\hspace*{12pt}(9) \textit{There is a map $\theta$ from $\{0, \dots, l-1\}$ onto $\{0, \dots, k-1\}$ with $\theta(0) = 0$, $\theta(l-1) = k-1$ and $w_q \in V(\sigma'(u_{\theta(q)}))$ for all $q$, such that for $0 \le q < l-1$, either $\theta(q+1) = \theta(q)$ and the vertex of $R^*$ between $w_q$ and $w_{q+1}$ is regular, or $\theta(q+1) = \theta(q)+1$ and the border edge of $\Gamma'_j$ between $w_q$ and $w_{q+1}$ is $\sigma'(e)$, where $e$ is the border edge of $\Gamma'_i$ between $u_{\theta(q)}$ and $u_{\theta(q)+1}$.}

Let $\theta(0) = 0$; then $w_0 \in V(\sigma'(u_0))$ since $w_0 = u_0 \in V(\Phi')$. Let $\theta(q) = p$ where $q < l-1$, and let $f$ be the border edge of $\Gamma'_j$ between $w_q$ and $w_{q+1}$. If $f$ is a short side, then $f$ is also a border edge of $\Gamma'_i$ and $\sigma'(f) = f$; and $w_q = u_p$, since $w_q \in V(\sigma'(w_q)) \cap V(\sigma'(u_p))$, and so $w_{q+1} = u_{p+1}$. If $f = \sigma'(e)$ where $e$ is not a short side, then $e$ borders $bd(W)$ since $\sigma'$ is linear, and $e, f$ are directed the same way since they border the same side of $\Phi'$; so $e$ is the border edge between $u_p$ and $u_{p+1}$, by the third condition in the definition of an inflation and since the sets $\sigma'(v)$ are disjoint. In these two cases we set $\theta(q+1) = p+1$; here $p+1 < k$, because the border edge through $r_0^*$ is its own image and is not $f$. Otherwise, by the second condition in the definition of ``linear'', $f$ is an edge of $\sigma'(v)$ for some $v$, and $v = u_p$ since $w_q \in V(\sigma'(u_p))$; then $f$ is not a short side, since short sides are their own images, so the vertex of $R^*$ in $f$ is regular, and we set $\theta(q+1) = p$. Since $w_{l-1} = u_{k-1}$ is in $\sigma'(u_{k-1})$ and in $\sigma'(u_{\theta(l-1)})$, we have $\theta(l-1) = k-1$; and so $\theta$ is onto, since it increases by at most one at each step. This proves (9).

\hspace*{12pt}(10) \textit{For $0 \le q < l$, the spokes $S_{w_q}$ and $S_{u_{\theta(q)}}$ have the same number of levels, and $h_j(w_q) = h_i(u_{\theta(q)})$.}

We use induction on $q$. For $q = 0$ the two spokes are equal, since $w_0 = u_0 \in N(W)$. Let $q < l-1$ and $\theta(q) = p$, and assume (10) holds for $q$. If the border edge between $w_q, w_{q+1}$ is a short side, then $\theta(q+1) = p+1$ and $w_{q+1} = u_{p+1} \in N(W)$ by the proof of (9), and the two spokes are equal. If $\theta(q+1) = p$, then the vertex of $R^*$ between $w_q$ and $w_{q+1}$ is regular, by (9). So $S_{w_{q+1}}$ and $S_{w_q}$ have the same inner end, by the definition of ``regular'' if $|V(H(W))| > 1$, and since every spoke ends in the only vertex of $H(W)$ otherwise. Hence they have the same number of levels and the same hub index, as we saw above, and (10) holds for $q+1$ since it holds for $q$. In the remaining case, by (9), $\theta(q+1) = p+1$ and the border edge between $w_q, w_{q+1}$ is $\sigma'(e)$, where $e$ is the border edge between $u_p, u_{p+1}$, and $e$ is not a short side. Let $u_{p+1}$ be the $k$th term of $\gamma_{\Gamma'_i}(e)$. By the third condition in the definition of an inflation, the $k$th term of $\gamma_{\Gamma'_j}(\sigma'(e))$ is in $V(\sigma'(u_{p+1}))$; it is $w_q$ or $w_{q+1}$, and it is not $w_q$, since $w_q \in V(\sigma'(u_p))$ and the sets $\sigma'(v)$ are disjoint. So $w_{q+1}$ is the $k$th term of $\gamma_{\Gamma'_j}(\sigma'(e))$. Now $\phi'_i(e) \le \phi'_j(\sigma'(e))$, and in both colours the march of the tail comes first; so by the definition of the order on $R$, there is a bijection mapping $\mu_i(u_{p+1})$ to $\mu_j(w_{q+1})$, and $h_i(u_{p+1}) = h_j(w_{q+1})$. This proves (10).

\hspace*{12pt}(11) \textit{Let $p = \theta(q)$, where either $q < l-1$ and $\theta(q+1) = p+1$, or $q = l-1$ (and then we put $w_l = w_0$, $u_k = u_0$). Then for all $t$, $w_q^t = w_{q+1}^t$ if and only if $u_p^t = u_{p+1}^t$.}

If the border edge between $w_q, w_{q+1}$ is a short side (as it is if $q = l-1$), then $w_q = u_p$ and $w_{q+1} = u_{p+1}$ are in $N(W)$, and their spokes are the same for $\Gamma_i$ and $\Gamma_j$. Otherwise it is $\sigma'(e)$, where $e$ is the border edge between $u_p, u_{p+1}$, and since $\phi'_i(e) \le \phi'_j(\sigma'(e))$ there is a bijection mapping $\mu_i(u_p)$ to $\mu_j(w_q)$ and $\mu_i(u_{p+1})$ to $\mu_j(w_{q+1})$, as in the proof of (10). This proves (11).

\hspace*{12pt}(12) \textit{If a vertex is at level $t$ of $S_{w_q}$ and at level $t'$ of $S_{w_{q'}}$, then $t = t'$ and $u_{\theta(q)}^t = u_{\theta(q')}^t$.}

Let $y$ be the vertex. By 13.2 the two spokes have the same inner end and the same final segment from $y$, and they have the same length; so $t = t'$. We may assume that $q < q'$. The segments of $S_{w_q}$ and $S_{w_{q'}}$ from their first vertices to $y$ form a path of $\Gamma^*$ in $W$ which meets $bd(W)$ only in its ends, by 13.1, and so it divides $W$ into two closed discs, each bounded by this path and one of the two arcs of $bd(W)$ between $w_q$ and $w_{q'}$. Since a spoke meets $V(H(W))$ only in its inner end, $U(H(W))$ meets this path at most in $y$, and so one of the two discs, $D$ say, contains no point of $U(H(W))$ except perhaps $y$. Let $w_s$ be a vertex of $\Gamma_j$ in the arc of $bd(W)$ in $D$. The spoke $S_{w_s}$ starts in $D$, ends in $V(H(W))$, and meets $bd(W)$ only in $w_s$; so it meets the path above in a vertex, and from its first such vertex on it coincides with $S_{w_q}$ or $S_{w_{q'}}$, by 13.2. Hence $y$ is a vertex of $S_{w_s}$, and at level $t$ of it, by the first part of the argument. Thus $y = w_s^t$ for all $s$ in one of the two cyclic intervals from $q$ to $q'$, namely $q, q+1, \dots, q'$ or $q', q'+1, \dots, l-1, 0, \dots, q$. Along this interval, at each step from $s$ to $s+1$ (or from $l-1$ to $0$), either $\theta(s+1) = \theta(s)$, or (11) applies and gives $u^t_{\theta(s)} = u^t_{\theta(s+1)}$, since $w_s^t = w_{s+1}^t$. So $u_{\theta(q)}^t = u_{\theta(q')}^t$. This proves (12).

Each vertex of $M(\Gamma_i, \mathscr{L}_i)$ in $W$ is on a spoke or on $H(W)$. For $W$ is the union of the discs $\Delta_{r^*}$ $(r^* \in R^* \cap bd(W))$ and of $ins(H(W))$ if $|V(H(W))| > 1$; and a vertex $x$ of $\Gamma_i$ in one of these discs, not on a spoke and not on $H(W)$, belongs to $V(A^-) \setminus \bar{\pi}(A)$ for the corresponding sector $A$, and so $x \notin V(A^c)$ and $x \notin V(M(\Gamma_i, \mathscr{L}_i))$. For $x \in V(M(\Gamma_i, \mathscr{L}_i))$, we say that $y \in V(\Gamma_j)$ is \textit{sent back} to $x$ if either $x \in V(\Gamma'_i)$ and $y \in V(\sigma'(x))$; or $y = w_q^t$ and $x = u_{\theta(q)}^t$ for some $W, q, t$; or $\theta(q+1) = \theta(q)$, $x = u_{\theta(q)}^t$, and $y$ is a vertex of the path $P_t$ of 13.7 between $w_q^t$ and $w_{q+1}^t$, for some $W, q, t$ (13.7 applies, since $S_{w_q}, S_{w_{q+1}}$ have the same inner end by (9)); or $y = x \in V(H(W))$ for some $W$. Let $\sigma(x)$ be the subgraph of $sk(\Gamma_j)$ induced on the vertices sent back to $x$, and let $\sigma(e) = \sigma'(e)$ for $e \in E(M(\Gamma_i, \mathscr{L}_i))$.

Each $\sigma(x)$ is non-null, by (9) and (10), and since $V(\Gamma) \cap V(H(W))$ is the same for all members of $\mathscr{D}'$. The sets $\sigma(x)$ are disjoint. For let $y$ be sent back to $x$. If $y$ is in $\Sigma'$ and in no $bd(W)$, then $y \in V(\sigma'(x))$, and $x$ is determined by $y$ since the sets $\sigma'(v)$ are disjoint. If $y$ is at level $t$ of $S_{w_q}$ (for instance $y = w_q$, with $t = 0$), then $x = u^t_{\theta(q)}$: by (9) for the first rule, by (12) for the second, also for the third since $y$ is then an end of its path, and as follows for the fourth. If $y$ is sent back to $x$ by the fourth rule, then $y = x \in V(H(W))$; so $y$ is the inner end of $S_{w_q}$, since a spoke meets $V(H(W))$ only in its inner end, and $t$ is the last level of $S_{w_q}$, and $h_j(w_q)$ is the position of $y$ in $\pi(H(W))$. By (10), $S_{u_{\theta(q)}}$ also has last level $t$, and $h_i(u_{\theta(q)}) = h_j(w_q)$; so $u^t_{\theta(q)} = y = x$. If $y$ is on $H(W)$ and on no spoke, then $x = y$. Otherwise $y$ is inside some disc $\Delta_{s^*}$ and on some path $P_t$, and these determine $x$, since the discs $\Delta_{s^*}$ have disjoint interiors and the paths of 13.7 are disjoint. Each $\sigma(x)$ is connected. For let $w_q^t$ and $w_{q'}^t$ be sent back to $x$ by the second rule. By the argument of (12) applied to $\Gamma_i$, $x$ is at level $t$ of $S_{u_p}$ for all $p$ in one of the two cyclic intervals from $\theta(q)$ to $\theta(q')$; so $w_s^t$ is sent back to $x$ for all $s$ in the corresponding interval from $q$ to $q'$, and $w_s^t, w_{s+1}^t$ are equal by (11) if $\theta(s+1) \ne \theta(s)$, and joined by $P_t \subseteq \sigma(x)$ if $\theta(s+1) = \theta(s)$. A vertex sent back to $x$ by the third rule is joined by its path to an end, which is sent back to $x$ by the second rule. A vertex sent back to $x$ by the fourth rule is $x$ itself. If $x = u^t_{\theta(q)}$, then $x$ is the inner end of $S_{u_{\theta(q)}}$ and $t$ is its last level, and so $w_q^t = x$ by (10), in the same way; and if $x$ is on no spoke, then $x$ is the only vertex of $\sigma(x)$. Finally let $x \in V(\Gamma'_i)$. Then $\sigma'(x)$ is a connected subdrawing of $sk(\Gamma'_j)$, and its edges not in $bd(\Sigma')$ are edges of $sk(\Gamma_j)$. If $w_q \in V(\sigma'(x))$ for some $W$ and $q$, then $x = u_{\theta(q)}$, by (9) and since the sets $\sigma'(v)$ are disjoint; so $w_q$ is sent back to $x = u^0_{\theta(q)}$ also by the second rule. Each edge of $\sigma'(x)$ in $bd(W)$ joins some $w_q, w_{q+1}$; it is not in $\sigma'(E(\Gamma'_i))$, since $\sigma'(x)$ contains at most one end of each such edge, and so $\theta(q+1) = \theta(q)$ by the proof of (9), and $w_q, w_{q+1}$ are joined in $\sigma(x)$ by $P_0$. Hence $\sigma(x)$ is connected. (If $x$ is in no $bd(W)$, then $\sigma(x) = \sigma'(x)$, and $\sigma'(x)$ has no edge in $bd(\Sigma')$.) Each edge $e$ of $M(\Gamma_i, \mathscr{L}_i)$ is disjoint from every $W$, so it is an internal edge of $\Gamma'_i$; hence $\sigma'(e)$ is an edge of $\Gamma_j$ disjoint from every $W$, that is, an edge of $M(\Gamma_j, \mathscr{L}_j)$, with $\phi_i(e) \le \phi_j(\sigma'(e))$, and its $h$th end is in $\sigma'(v) \subseteq \sigma(v)$ where $v$ is the $h$th end of $e$. So $\sigma$ is an inflation, and the first and last conditions in the definition of an outline hold.

For the others, let $A \in \mathscr{L}_i$ be the sector at the border edge between $u_p, u_{p+1}$. By (9) and the definition of $\tau$, $\tau(A)$ is the sector at the border edge between $w_q, w_{q+1}$, where $\theta(q) = p$ and $\theta(q+1) = p+1$, or $q = l-1$ and $p = k-1$; so $\tau$ is one-to-one. By the bijection of (11), or because $\pi(r^*)$ is fixed if $r^* \in S(W)$, $|\bar{\pi}(\tau(A))| = |\bar{\pi}(A)|$, and the $n$th terms of $\pi(A)$ and $\pi(\tau(A))$ are at the same level of $S_{u_p}$ and $S_{w_q}$, or of $S_{u_{p+1}}$ and $S_{w_{q+1}}$; so the second is sent back to the first. The hub sectors have the same march $\pi(H(W))$. Finally, the vertices of $\tau(A)$ sent back to $x$ are on the two spokes of $\tau(A)$ or on $H(W)$, since the vertices of $\tau(A)$ in $\Sigma'$ are $w_q, w_{q+1}$, and the paths $P_t$ inside other discs do not meet $\tau(A)$; and two such vertices sent back to the same $x$ are equal, by 13.2, (10) and (11). The same holds for the hub sector, since a vertex of $\Gamma_j$ inside $H(W)$ that is sent back to some $x$ is on $H(W)$, and equals $x$. So $|V(\sigma(x)) \cap V(\tau(A))| \le 1$ for all $x$.}\fixnote{fix:14.1-outline-proof-added}{The paper omits this verification, and refers to the proof of Theorem 9.1 of [11]. The text above replaces the construction of $\sigma$ by one in which $\sigma(x)$ is read off along the map $\theta$ of each wheel, and gives the verification. This is the argument used in the formalization.}

Thus $\tau \colon \mathscr{L}_i \to \mathscr{L}_j$ is an outline and by (8) the proof is complete. $\square$

\sectionhead{15. Conclusion}

The proof will be completed by using one further result. To prove it we need the following lemma.

15.1. \textit{Let $\Gamma$ be a 2-cell drawing with $E(\Gamma) \ne \emptyset$ in a surface $\Sigma$ with $\text{bd}(\Sigma) = \emptyset$, and let $\mathscr{T}$ be a tangle in $\Gamma$ with $\text{rep}(\mathscr{T}) \ge \theta$. Let $r$ be a region of $\Gamma$, and let $X \subseteq V(\Gamma) \cap \bar{r}$ with $|X| \le \theta$. Then the following are equivalent:}

\hspace*{12pt}$\bullet$ \textit{there exists $(A, B) \in \mathscr{T}$ of order $< |X|$ with $X \subseteq V(A)$,}

\hspace*{12pt}$\bullet$ \textit{there is a circuit $C$ of $\Gamma^*$ of length $< 2\theta$, such that $|X \cap \text{ins}(C)| > \frac{1}{2}|E(C)|$ and $r^* \in \text{ins}(C)$.}

\textbf{Proof.} That the second statement implies the first is easy, for if $C$ is as in the second statement, let
\begin{align*}
V(A) &= V(\Gamma \cap \text{ins}(C)) \cup X \\
E(A) &= E(\Gamma \cap \text{ins}(C)) \\
B &= \Gamma \cap \text{out}(C).
\end{align*}
Then $(A, B) \in \mathscr{T}$, $(A, B)$ has order $\frac{1}{2}|E(C)| + |X \setminus \text{ins}(C)| < |X|$, and $X \subseteq V(A)$; and so the first statement holds.

For the converse, let $(A, B)$ satisfy the first statement. Let $G \subseteq \Gamma^*$ be the subdrawing consisting of all edges of $\Gamma^*$ split by $(A, B)$, and their ends. Let $C_1, \dots, C_k$ be all the circuits $C$ of $G$ with $\text{ins}(C)$ maximal. By Theorem 6.3 of [6], every edge of $A$
